\section{从 SSL 到 BiT：规模化预训练需要整套 recipe}

定义测量后，讲座回顾三条获得视觉表示的路线：self-supervised pretraining、semi-supervised learning 与 large-scale supervised pretraining。重点不是给三者排永久名次，而是展示视觉领域的性能取决于 labels、data scale、model class、optimization 与 evaluation 共同作用，不能简单复制语言模型的单一成功叙事。

\subsection{路线图：五个问题而不是一个答案}

概览 slide 将 SSL、半监督、BiT、ViT 与 further scaling 放在同一序列中。读者应把它理解为不断减少人工先验、扩大数据与模型、再用 transfer benchmark 复核的实验程序。每一站都同时改变监督信号、模型类或训练规模，因此本节先建立变量账本，再讨论结果，避免把完整系统增益错误归给一个新名词。

\lecturefigure{slide-10-overview.jpg}{讲座路线从 SSL、半监督与 BiT 走向 ViT、Scaling ViT 和 MLP-Mixer。}{官方视频 07:20。}

\begin{knowledgebox}{读图：每一步改变了哪些变量}
SSL 改变监督信号；半监督混合 labeled/unlabeled data；BiT 主要扩大监督数据与模型；ViT 改变 architecture prior；Scaling ViT 改变容量、shape 与 schedule；Mixer 进一步移除 attention。若多个变量同时变化，结论只能归因于完整 recipe，除非有消融实验隔离贡献。
\end{knowledgebox}

\subsection{Self-supervised learning：视觉代理任务不自动等于语义}

\term{self-supervised learning}（自监督学习）从数据本身构造监督信号，例如预测遮挡区域、对比不同增强视图或重建输入。它减少人工标签依赖，但 proxy objective 可能奖励纹理、增强不变性或数据集捷径，而非所有 downstream 所需结构。

\lecturefigure{slide-11-self-supervised-review.jpg}{早期视觉 SSL 结果说明架构与代理任务选择会显著影响迁移。}{官方视频 07:56。}

\begin{knowledgebox}{读图：不要只比较柱子高度}
图中方法同时改变 pretext task、architecture 与 training recipe。更高 mean score 可以说明某组合更有效，却不能单独证明“无标签数据优于标签”或“某代理任务发现了语义”。应检查同 backbone、同 compute、同 augmentation 的 controlled comparison，并观察 structured/specialized tasks 是否同步受益。
\end{knowledgebox}

\teachervoice{Lucas 提醒，NLP 中自然出现的 next-token target 在视觉里没有完全对应的唯一目标。课堂提示：无标签不等于无设计；研究者把大量先验放进 augmentation、positive-pair construction、masking pattern 与 architecture。}

\subsection{Semi-supervised learning：标签价值取决于模型类}

\term{semi-supervised learning}（半监督学习）同时使用少量 labeled data 与大量 unlabeled data，常见机制包括 pseudo-labeling、consistency regularization 与 teacher-student training。它看似是折中方案，但模型类别和标签预算会改变收益：新的 architecture 可能让相同标签产生更大增益，也可能让伪标签错误被反复放大。

\lecturefigure{slide-12-semi-supervised-review.jpg}{半监督路线比较标签预算、模型类与无标签数据的组合。}{官方视频 08:36。}

\begin{warningbox}{Pseudo-label feedback loop}
若 teacher 在某类样本上系统性错误，pseudo-label 会把错误变成训练目标。需要 confidence threshold、class balance、strong augmentation、held-out calibration 与 subgroup analysis。总体准确率提升不能证明少数 domain 或长尾类别没有退化。
\end{warningbox}

\subsection{Big Transfer：耐心、容量与数据必须一起扩展}

BiT 使用大规模监督预训练的 ResNet 系列研究 transfer。其持久价值不只是一个 checkpoint，而是给出工程规律：更多数据需要更大容量；大模型可能在训练早期看起来更差；下游适配应根据数据规模选择 schedule 与 regularization；鲁棒性和去重必须成为报告的一部分。

\lecturefigure{slide-13-bit-summary.jpg}{BiT 总结了耐心训练、同时扩大数据与模型、few-shot transfer 和 robustness。}{官方视频 23:16；Kolesnikov et al. 2019。}

\begin{knowledgebox}{读图：三条经验如何连接}
左侧训练曲线提示大模型需要更长优化才能显现优势；中间散点说明数据规模与模型容量共同决定 transfer；右侧 few-shot 曲线显示表示可降低下游样本需求；下方 ObjectNet 结果检验分布变化。四部分共同说明：upstream loss、in-distribution accuracy、few-shot transfer 与 robustness 是不同证据，不能只保留最好看的一个。
\end{knowledgebox}

\begin{importantbox}{容量匹配的直觉}
设数据规模为 $D$、模型容量为 $N$、训练计算为 $C$。若固定很小的 $N$ 而只增大 $D$，模型可能无法吸收新增结构；若固定很小的 $D$ 而增大 $N$，则更易过拟合或依赖强 regularization。实际目标是寻找受 compute 约束的组合
\[
(N^{\star},D^{\star})=\arg\min_{N,D:\,C(N,D)\le C_0}\mathcal{L}_{\text{transfer}}(N,D),
\]
而不是分别最大化参数或样本数。
\end{importantbox}

\lecturefigure{slide-14-bit-large-scale-pretraining.jpg}{BiT 在大规模预训练与多任务迁移之间建立经验联系。}{官方视频 23:44。}

\begin{knowledgebox}{读图：上游规模不是最终目标}
横轴对应时间或规模演进，重点是大规模 supervised pretraining 产生可复用 checkpoint。图上领先点不表示所有任务都同幅度提升；应回到 VTAB/few-shot/ObjectNet 分组结果，确认收益是否跨任务、是否超过额外 compute 成本，并检查测试分布是否与预训练重叠。
\end{knowledgebox}

\teachervoice{老师反复强调“be patient”：更深更大的模型在早期训练可能落后，不能用短跑实验否定规模化方案。实践经验是先确保 schedule 足够长、optimization 稳定，再比较最终 transfer；否则是在比较训练成熟度，而非模型潜力。}

\subsection{Deduplication：没有数据审计就没有可信迁移}

大规模数据集可能包含 downstream test images 的相同或近重复版本。\term{deduplication}（去重）不仅移除完全相同的文件，还要检测裁剪、压缩、缩放、颜色变化与轻度编辑后的 near-duplicates。若测试图进入 pretraining，few-shot transfer 会被记忆伪装成泛化。

\lecturefigure{slide-15-deduplication.jpg}{BiT 对多个下游数据集执行近重复检测，展示潜在泄漏样本。}{官方视频 25:00。}

\begin{warningbox}{数据泄漏的四个层级}
\begin{enumerate}
  \item \textbf{Exact duplicate}：文件或像素完全相同；
  \item \textbf{Near duplicate}：裁剪、缩放、压缩或水印变化；
  \item \textbf{Semantic duplicate}：同一对象/场景的相邻照片；
  \item \textbf{Label leakage}：文件名、网页文字或 metadata 暴露答案。
\end{enumerate}
仅靠哈希可处理第一层，后续层级需要 perceptual embedding、nearest-neighbor review 与 provenance records。
\end{warningbox}

\subsection{本章小结}

视觉预训练没有单一魔法目标。SSL 把先验移入代理任务，半监督需要控制伪标签反馈，BiT 证明数据、容量、训练长度与迁移协议必须一起扩展。任何大规模 transfer 结论都应附带去重和分布审计。

